% =========================================================================
% ENCODAGE ET LANGUE
% =========================================================================
\usepackage[T1]{fontenc}
\usepackage[french]{babel}
\usepackage{lmodern}

% =========================================================================
% GEOMETRIE ET EN-TETES
% =========================================================================
\usepackage[a4paper, margin=2.5cm, headheight=14pt]{geometry}
\usepackage{fancyhdr}
\pagestyle{fancy}
\fancyhf{}
\fancyhead[L]{\small\scshape Analyse III}
\fancyhead[R]{\small\leftmark}
\fancyfoot[C]{\thepage}
\renewcommand{\headrulewidth}{0.4pt}

% =========================================================================
% PACKAGES MATHEMATIQUES
% =========================================================================
\usepackage{amsmath, amsfonts, amssymb, amsthm}
\usepackage{mathtools}
\usepackage{physics} % Pour \pdv, \grad, \div, \curl, \bra, \ket, etc.
\usepackage{bm}      % Symboles gras en mode mathématique

% Ens. de nombres usuels
\newcommand{\R}{\mathbb{R}}
\newcommand{\C}{\mathbb{C}}
\newcommand{\N}{\mathbb{N}}
\newcommand{\Z}{\mathbb{Z}}
\newcommand{\K}{\mathbb{K}}

% Operateurs d'analyse
\DeclareMathOperator{\sinc}{sinc}
\DeclareMathOperator{\supp}{supp}
\DeclareMathOperator{\res}{Res}

% =========================================================================
% TIKZ ET GRAPHISME
% =========================================================================
\usepackage{tikz}
\usepackage{pgfplots}
\pgfplotsset{compat=1.18}
\usetikzlibrary{arrows.meta, calc, positioning, shapes, patterns, decorations.pathmorphing}

% =========================================================================
% COULEURS & TCOLORBOX
% =========================================================================
\usepackage[x11names, dvipsnames, table]{xcolor}
\usepackage[most]{tcolorbox}

% Definition des teintes (Palettes professionnelles et contrastees)
\definecolor{ThmColor}{HTML}{1A365D}   % Bleu Nuit
\definecolor{PropColor}{HTML}{2B6CB0}  % Bleu Petrole
\definecolor{DefColor}{HTML}{1C4532}   % Vert Foret
\definecolor{ExColor}{HTML}{744210}    % Ambre / Brun
\definecolor{RemColor}{HTML}{4A5568}   % Slate / Anthracite
\definecolor{ProofColor}{HTML}{4A154B} % Violet Sombre

% Style de base pour les boites
\tcbset{
    basebox/.style={
        enhanced,
        breakable,
        boxrule=0pt,
        frame hidden,
        arc=2pt,
        outer arc=2pt,
        top=8pt,
        bottom=8pt,
        left=10pt,
        right=10pt,
        separator sign={. \ },
        fonttitle=\bfseries,
    }
}

% Env 1 : Theoreme
\newtcbtheorem[number within=section]{theoreme}{Théorème}{
    basebox,
    borderline west={3.5pt}{0pt}{ThmColor},
    colback=ThmColor!5!white,
    coltitle=ThmColor,
    colbacktitle=ThmColor!5!white,
    attach title to upper,
}{thm}

% Env 2 : Proposition
\newtcbtheorem[use counter from=theoreme]{proposition}{Proposition}{
    basebox,
    borderline west={3.5pt}{0pt}{PropColor},
    colback=PropColor!5!white,
    coltitle=PropColor,
    colbacktitle=PropColor!5!white,
    attach title to upper,
}{prop}

% Env 3 : Definition
\newtcbtheorem[use counter from=theoreme]{definition}{Définition}{
    basebox,
    borderline west={3.5pt}{0pt}{DefColor},
    colback=DefColor!5!white,
    coltitle=DefColor,
    colbacktitle=DefColor!5!white,
    attach title to upper,
}{def}

% Env 4 : Exemple (Sans numérotation obligatoire)
\newtcolorbox{exemple}[1][]{
    basebox,
    borderline west={3.5pt}{0pt}{ExColor},
    colback=ExColor!5!white,
    coltitle=ExColor,
    title=\textbf{Exemple\ifx\relax#1\relax\else\ (#1)\fi .},
    attach title to upper,
}

% Env 5 : Remarque
\newtcolorbox{remarque}[1][]{
    basebox,
    borderline west={3.5pt}{0pt}{RemColor},
    colback=RemColor!5!white,
    coltitle=RemColor,
    title=\textbf{Remarque\ifx\relax#1\relax\else\ (#1)\fi .},
    attach title to upper,
}

% Env 6 : Demonstration
\newtcolorbox{demonstration}[1][]{
    basebox,
    borderline west={2pt}{0pt}{ProofColor},
    colback=ProofColor!3!white,
    coltitle=ProofColor,
    title=\textbf{Démonstration\ifx\relax#1\relax\else\ (#1)\fi .},
    attach title to upper,
    after upper={\hfill $\square$}
}

% =========================================================================
% HYPERREF & OUTILS DIVERS
% =========================================================================
\usepackage{microtype}
\usepackage{hyperref}
\hypersetup{
    colorlinks=true,
    linkcolor=ThmColor,
    citecolor=DefColor,
    urlcolor=PropColor,
    pdfauthor={Étudiant EPFL},
    pdftitle={Polycopié Analyse III}
}